\section{Introduction}
\label{sec:intro}

Writing test suites manually is time-consuming and error-prone, consuming up to 50\% of total software development budgets~\cite{myers2011art, evosuite, randoop, swebench}. In repository-level software maintenance, comprehensive unit testing is essential to prevent regressions and validate intended system semantics~\cite{docprompting, repobench, cse}. Consequently, automated test generation has emerged as a fundamental research area in automated software engineering, aiming to reduce developer burden while maintaining high code quality and software reliability.

Recent studies have shown that Large Language Models (LLMs) can synthesize readable unit tests and achieve high code coverage across diverse programming languages~\cite{bareiss2022code, lemieux2023codamosa, xie2023chatunitest}. Advanced generative approaches leverage deep pre-trained semantic knowledge to generate human-like assertion statements and contextual setup code, outperforming traditional search-based software testing (SBST) techniques in terms of semantic readability and intent alignment~\cite{siddiq2023exploratory, schafer2023empirical}.

However, no study has evaluated interactive terminal execution feedback for repository-level test suite generation on documentation-derived intent benchmarks using real-time execution feedback loops. Existing LLM-based test generation tools operate predominantly in a static Plain Context paradigm, relying on single-pass code completion without dynamic execution feedback~\cite{testexplora}. Lacking runtime execution signals, these static approaches suffer severely from \textit{compliance bias}: models passively generate test assertions that conform to existing (defective) implementations, causing generated tests to pass on buggy codebases and fail to detect latent logic bugs~\cite{testexplora, sweagent}.

In this paper, we investigate the effectiveness of interactive terminal execution feedback in breaking compliance bias for repository-level test suite generation. We propose an open-source multi-agent framework featuring an Agentic Exploration Agent (DeepSeek-V3 operating via an interactive SWE-agent CLI bash terminal) integrated with a Code Action Fixer Agent (Llama-3.3-70B). By executing test suites in real-time, observing \texttt{pytest} tracebacks, and iteratively refining adversarial assertions, our system proactively isolates repository-level defects.

To summarize, this paper contributes:
\begin{enumerate}
    \item The first empirical evaluation of interactive terminal execution feedback for proactive repository-level test suite generation using Fail-to-Pass (F2P) rate on the TestExplora benchmark.
    \item Experimental results on $N=1,552$ tasks ($4,656$ agent runs across 16 parallel shards) showing that our framework achieves an average Pass@1 F2P rate of 35.05\% (more than double the 16.60\% static baseline) and a Pass@3 rate of 71.59\% ($p = 0.000000$, $\text{Cliff's } \delta = 0.3018$).
    \item A publicly available replication package including full execution logs, scripts, merged dataset, and figures at \url{https://github.com/lctx9/group4}.
\end{enumerate}

The rest of this paper is structured as follows. Section~\ref{sec:related} reviews related work in automated test generation and compliance bias. Section~\ref{sec:method} details our multi-agent architecture and experimental methodology. Section~\ref{sec:results} presents our empirical results and statistical analysis. Section~\ref{sec:discussion} provides qualitative discussion and failure mode analysis. Section~\ref{sec:threats} addresses threats to validity, and Section~\ref{sec:conclusion} concludes the paper and outlines future directions.
